\chapter{Materials selection}
\label{vol05:ch10}

\epigraph{Materials are selected by performance, by cost, and by the
engineer's willingness to be honest about which performance the
application actually needs. The Ashby method is the discipline of
making that honesty visible on a chart.}{Adapted from M.~F. Ashby,
working scope of materials selection in mechanical design \cite[ch.~1,
pp.~1--13]{text:ashby2017}.}

\chapmeta{Half-life: durable. Archetypes: optimisation, ethics. Exercise target: 28. Project track: design study with property-chart analysis. Hazard class: none. Reader path default: standard, with sections explicitly tagged. Named cases: Hubble primary mirror (1990) referenced for materials-process trade-off; price-only-selection failures described by mechanism only.}

\noindent
A materials-selection decision is a discrete choice from a list of
candidates against a set of constraints and a chosen objective. The
chapter develops the Ashby method that organises the choice: property
charts that locate the candidates in property space, performance
indices that rank candidates against a chosen objective, constraints
that reject infeasible candidates, multi-objective trade-offs when the
indices disagree, and the non-mechanical constraints (cost,
manufacturability, availability, sustainability) that bind real
engineering decisions. The closing section names the canonical failure
mode of price-only selection. The chapter is the synthesis chapter of
the volume: every preceding chapter contributes properties; selection
is the discipline of using them.

\section{The Ashby method introduced}\pathtag{core}

The materials-selection problem has the form: choose a material
$M$ from a candidate set $\mathcal{M}$ that satisfies a set of
constraints $\{g_i(M) \leq 0\}$ (or $\geq 0$) and minimises (or
maximises) an objective $f(M)$. The candidate set is the set of all
engineering materials; the constraints are the engineering
requirements; the objective is the engineering goal. The Ashby method
\cite[ch.~3]{text:ashby2017} organises the problem in four steps.

\begin{enumerate}[leftmargin=*]
\item \engterm{Translation}. Convert the design requirements into the
  language of materials and processes: function, constraints,
  objective, free variables. Function: what does the component do?
  Constraints: what must it do without fail? Objective: what should
  it do best? Free variables: which dimensions or properties may the
  designer choose?
\item \engterm{Screening}. Apply the constraints to reject materials
  that cannot meet the requirements. The constraints are absolute:
  a material that fails any single constraint is rejected.
\item \engterm{Ranking}. Order the surviving materials by their
  value of a performance index that combines the relevant
  properties with the objective and the free variables. The
  performance index is the figure of merit; the highest (or lowest)
  index identifies the best material.
\item \engterm{Documentation}. Look up the candidates' working data,
  manufacturing routes, and supply chains. Confirm or refute the
  shortlist with the full data.
\end{enumerate}

\engterm{The function-constraints-objective-free variable
formulation}. The discipline of stating the problem in these four
fields is the working content of the method. A working example: a
light, stiff tie rod of fixed length $L$ that must carry a tensile
load $F$ at a stress below the material's yield strength. Function:
tie rod, tension. Constraints: length $L$ fixed; cross-section $A$
free; must not yield ($\sigma_y A \geq F$). Objective: minimise mass
$m = \rho A L$. Free variable: cross-section $A$.

Eliminating $A$ via the constraint, $A = F / \sigma_y$, the mass
becomes $m = F L \rho / \sigma_y$. The mass is minimised by
maximising $\sigma_y / \rho$, the \engterm{specific strength}. The
performance index for the light, strong tie rod is $M = \sigma_y /
\rho$. Materials with higher $M$ give lighter tie rods at the same
load.

\engterm{Indices for the common loading cases}. Ashby's catalogue
runs to hundreds of indices for the common combinations of geometry,
loading, and objective \cite[Appendix B]{text:ashby2017}. The
working core:

\begin{itemize}[leftmargin=*]
\item Light, stiff tie rod (tension, stiffness objective): $M =
  E / \rho$.
\item Light, strong tie rod (tension, strength objective): $M =
  \sigma_y / \rho$.
\item Light, stiff beam (bending, stiffness objective, free
  thickness): $M = E^{1/2} / \rho$.
\item Light, strong beam (bending, strength objective, free
  thickness): $M = \sigma_y^{2/3} / \rho$.
\item Light, stiff panel (bending, stiffness objective, free
  thickness): $M = E^{1/3} / \rho$.
\item Light, strong panel (bending, strength objective, free
  thickness): $M = \sigma_y^{1/2} / \rho$.
\item Cheap, stiff beam: $M = E^{1/2} / (\rho C_m)$ with $C_m$ the
  material cost per kilogram.
\item Pressure vessel (thin-walled, yield-limited): mass scales as
  $\rho R^2 P / \sigma_y$, so the index is $\sigma_y / \rho$ for
  mass-limited service and $\sigma_y$ alone for cost-limited service.
\end{itemize}

The structure is the same in every case: the geometry sets the
exponents in the index, the objective sets the numerator and
denominator, the free variable identifies the property to be
substituted. The exponents are not arbitrary: they follow from the
mechanics of the loading case. The reader who internalises the
derivation of three or four indices owns the working machinery.

\begin{principle}[The performance index makes the materials choice a
ranking problem]
A design requirement, properly formulated, reduces to maximising or
minimising a property combination over a candidate set. The property
combination is the performance index. The materials choice is the
material with the highest (or lowest) index that satisfies all the
constraints. The discipline of the method is making the index
visible: the design intent is in the index, not in the chosen
material.
\end{principle}

\section{Property charts}\pathtag{core}

A \engterm{property chart} plots one material property against
another for all engineering materials. The chart is a log-log plot
covering five to ten decades on each axis, with materials grouped
into families (metals, ceramics, polymers, composites, foams,
elastomers, woods). Each family occupies a bounded region of the
chart; specific materials are points or small clouds.

\engterm{The canonical Ashby charts}. The working charts \cite[
ch.~4]{text:ashby2017} include:

\begin{itemize}[leftmargin=*]
\item \engterm{Modulus-density chart} ($E$ versus $\rho$): the
  most-used chart. Metals cluster at $\rho = 1$ -- $20\,\text{g}/
  \text{cm}^3$, $E = 50$ -- $400\,\text{GPa}$. Ceramics at $\rho =
  2$ -- $5\,\text{g}/\text{cm}^3$, $E = 100$ -- $1{,}000\,\text{GPa}$
  (diamond at the top). Polymers at $\rho = 1$ -- $2\,\text{g}/
  \text{cm}^3$, $E = 1$ -- $5\,\text{GPa}$. Composites at $\rho =
  1.5$ -- $2.5\,\text{g}/\text{cm}^3$, $E = 30$ -- $200\,\text{GPa}$.
  Foams at $\rho = 0.05$ -- $0.5\,\text{g}/\text{cm}^3$, $E = 0.01$
  -- $1\,\text{GPa}$.
\item \engterm{Strength-density chart} ($\sigma_y$ or $\sigma_f$
  versus $\rho$): the comparative chart for tensile applications.
\item \engterm{Strength-toughness chart} ($\sigma_y$ versus $K_{Ic}$):
  the chart for damage-tolerant design.
\item \engterm{Modulus-strength chart} ($E$ versus $\sigma_y$): the
  chart for buckling-versus-yield trade-offs.
\item \engterm{Thermal conductivity vs.\ electrical conductivity}:
  separates heat-exchanger materials (high $k$, high $\sigma_e$:
  metals) from heat-shielding materials (low $k$, low $\sigma_e$:
  ceramics and refractories) and from heat-shielding electrical
  insulators (low $k$, high $1/\sigma_e$: ceramics, polymers).
\item \engterm{Cost-modulus} and \engterm{cost-strength} charts: the
  charts for cost-bound design.
\end{itemize}

\engterm{Index lines on a chart}. A performance index $M = X^a /
Y^b$ corresponds to a line of constant index on a log-log chart with
$X$ and $Y$ as axes:
\[
\log Y = \frac{a}{b} \log X - \frac{1}{b} \log M.
\]
The slope of the constant-index line on the $\log Y$ versus $\log X$
chart is $a/b$. Materials above the line have higher index; materials
below have lower. The line is moved up across the chart until it just
touches the highest-index materials in the feasible set; those
materials are the shortlist.

\engterm{Working example: light stiff panel}. The performance index
for a light, stiff bending panel is $M = E^{1/3} / \rho$. On the
modulus-density chart with $E$ on the vertical axis and $\rho$ on the
horizontal, the constant-index line has slope $\log E / \log \rho =
3$. Sweeping a slope-three line from upper-left toward lower-right
identifies the panel candidates in order: technical foams (lowest
$\rho$, modest $E$, but $E^{1/3}/\rho$ still high), wood (along the
grain), aluminium-honeycomb sandwiches, magnesium alloys, balsa,
carbon-fibre composites, and titanium. For a panel that must also
support a point load (a constraint on yield strength), the screening
constraint rejects the foams; the shortlist becomes wood, balsa,
carbon-fibre, and magnesium.

\section{Performance indices}\pathtag{standard}

The construction of a performance index is the engineering content of
the method. The index follows mechanically from four ingredients: the
objective, the loading mode, the free variable, and the binding
constraint.

\engterm{Derivation in standard form}. The objective is the quantity
to be minimised or maximised: mass, cost, embodied energy, response
time. The objective is expressed in terms of the free variable and
the geometric quantities of the part: $m = \rho V$, $C = m C_m$ with
$C_m$ the material cost per kilogram, and so on. The constraint is
the mechanical or thermal requirement: a stiffness, a strength, a
critical temperature, a buckling load. The constraint relates the
free variable to the loading and the material properties: $E I/L^3 =
K$ for a stiff beam, $\sigma A = F$ for a strong tie rod, $E A/L = K$
for a stiff column. The free variable is eliminated by solving the
constraint; the resulting objective is a product of the loading and
geometric quantities (constant for the design problem) and a property
group (the index). The property group is the index.

\engterm{A two-page table of indices}. Ashby provides the canonical
table of indices for the common combinations of objective, loading,
and geometry \cite[Appendix B]{text:ashby2017}. The exponents follow
mechanically from the derivation:

\begin{itemize}[leftmargin=*]
\item Axial tie at given length, $\sigma_y$ binding: $M = \sigma_y /
  \rho$ for minimum mass; $M = \sigma_y / (\rho C_m)$ for minimum cost.
\item Beam at given length, $E$ binding, free thickness, free width:
  $M = E^{1/2} / \rho$ for minimum mass.
\item Beam at given length, $\sigma_y$ binding, free thickness, free
  width: $M = \sigma_y^{2/3} / \rho$ for minimum mass.
\item Plate at given dimensions, $E$ binding, free thickness: $M =
  E^{1/3} / \rho$ for minimum mass.
\item Column at given length, Euler buckling binding, free thickness,
  free width: $M = E^{1/2} / \rho$.
\item Heat-storage element, $\rho c_p T_{\max}$ binding, free volume:
  $M = \rho c_p T_{\max}$.
\item Heat-exchange tube, $k$ binding, free wall thickness: $M = k$.
\item Spring, $\sigma_y^2 / E$ binding: $M = \sigma_y^2 / (E \rho)$
  for minimum mass; $M = \sigma_y^2 / E$ for minimum volume.
\end{itemize}

The exponents differ from the tension case because the binding
constraint differs. A beam's bending stiffness scales as $I = b d^3/12$,
and the free variables (width $b$, depth $d$, or thickness $t$) absorb
different powers of $E$ when solving the constraint. The result is the
specific-stiffness exponents $1$, $1/2$, $1/3$ for tension,
free-width bending, free-thickness plate bending. The reader who
derives the three by hand can produce the rest of the catalogue in
minutes.

\engterm{The index for damage-tolerant design}. A part that must
tolerate a crack of length $a$ without fast fracture has the
constraint $K_I = Y \sigma \sqrt{\pi a} \leq K_{Ic}$. The working
stress is $\sigma = K_{Ic} / (Y \sqrt{\pi a})$, and the working load
in tension is $F = \sigma A$. Eliminating $A$ from the mass
equation, the index for a damage-tolerant light tie is $K_{Ic} /
\rho$, and for a damage-tolerant light beam (free thickness) is
$K_{Ic}^{2/3} / \rho$. The damage-tolerant indices put fracture-
toughness on the numerator; their constant-index lines have steep
slopes on a modulus-strength chart and pick materials with high
$K_{Ic}$, even at a strength penalty.

\section{Constraints and objectives}\pathtag{standard}

Real design problems impose multiple constraints simultaneously; the
discipline of the method is naming each one explicitly and applying
it as a screening cut before ranking begins.

\engterm{Typical constraint families}.

\begin{itemize}[leftmargin=*]
\item \engterm{Mechanical}: stiffness, strength, fracture toughness,
  fatigue endurance, creep at the operating temperature, buckling
  load.
\item \engterm{Thermal}: maximum service temperature, minimum service
  temperature, thermal-shock resistance, thermal expansion (matching
  to mating parts), thermal conductivity.
\item \engterm{Electrical and electromagnetic}: conductivity,
  dielectric strength, dielectric loss, magnetic permeability.
\item \engterm{Chemical and environmental}: corrosion resistance in
  the service environment, water absorption, UV stability,
  biocompatibility.
\item \engterm{Optical}: transparency, reflectivity, refractive index.
\item \engterm{Acoustic}: damping, sound speed.
\item \engterm{Process}: castability, formability, weldability,
  machinability, joinability, surface-finish capability.
\item \engterm{Regulatory}: RoHS and REACH compliance, food-contact
  approval, ATEX certification, biocompatibility (ISO 10993),
  regulatory restriction on lead, cadmium, hexavalent chromium.
\item \engterm{Supply chain}: availability in the required quantity,
  lead time, geographic source diversity.
\item \engterm{Other}: aesthetic, cultural, sustainability.
\end{itemize}

The screening step rejects every candidate that fails any constraint.
The shortlist after screening is the candidate set for ranking. A
poorly stated constraint admits candidates that should have been
rejected; a missing constraint admits candidates that fail in service.
The discipline of writing the constraint list explicitly is the
discipline of refusing to discover the missing constraint in the
field.

\engterm{The objective}. The objective is the single quantity the
design is meant to optimise. Common objectives:

\begin{itemize}[leftmargin=*]
\item Minimise mass (aerospace structure, sporting goods, hand tools,
  vehicles).
\item Minimise cost (consumer products, structural infrastructure,
  packaging).
\item Minimise embodied energy or carbon (sustainability-bound
  design).
\item Minimise volume (consumer-electronics housings, medical
  implants).
\item Maximise performance under a binding cost or mass limit.
\end{itemize}

A design with two objectives is not a single-objective problem; it is
a multi-objective problem treated in section~10.5.

\engterm{The free variable}. The free variable is what the designer
can change: typically a dimension (cross-section, thickness, length,
diameter). The choice of free variable determines the form of the
performance index. A beam with fixed cross-section is a different
problem from a beam with free thickness; the index changes
correspondingly.

\section{Multi-objective trade-offs}\pathtag{standard}

When the design has two or more objectives (light and cheap; strong
and damage-tolerant; conductive and corrosion-resistant), the
materials choice becomes a Pareto-optimisation problem. The materials
that lie on the trade-off front are the materials that cannot be
improved on one objective without degrading the other.

\engterm{The Pareto-front method}. Plot the materials on a chart with
the two objectives as axes (or transformed properties that align with
the objectives). The set of materials that lie on the lower-and-leftward
boundary (for objectives both minimised) form the Pareto front; the
materials inside are dominated and can be eliminated. The Pareto
front is the working shortlist.

The choice within the Pareto front requires a value statement: how
much of objective $A$ is the engineer willing to trade for objective
$B$? The trade is made explicit by a \engterm{utility function} of
the form $U = w_A M_A + w_B M_B$ with weights $w_A$ and $w_B$
reflecting the relative value of the two objectives in the
application. The utility-maximising material on the Pareto front is
the design's chosen material.

\engterm{Working example: light, cheap bicycle frame}. The objective
is to minimise both mass and cost for a bicycle frame of given
stiffness. On a chart of mass (normalised by stiffness$^{1/3}$ for a
free-thickness tube under bending) versus cost (normalised by the
same), the Pareto front runs from low-mass-high-cost (carbon-fibre
composite) through intermediate (titanium, aluminium 7075) to
high-mass-low-cost (mild steel, wood). The chosen material depends on
whether the bike is a road racer (mass dominates utility), a
commuter (cost dominates), or a touring frame (both matter,
intermediate).

\engterm{Sensitivity to the weights}. A Pareto-front analysis is
robust against small changes in the utility weights: the same
material remains optimal across a wide range of plausible weights.
The analysis is sensitive at the kinks of the front: at a kink, a
small change in weight shifts the optimum to a different material
family. The kinks are where the design's character changes.

\section{Cost, manufacturability, availability, sustainability as constraints}\pathtag{standard}

A material's selection in real engineering is bound by four
non-mechanical considerations the chapter has so far treated only
implicitly.

\engterm{Cost}. The material's cost per kilogram, current as of the
design year, and projected through the manufacturing campaign. Working
ranges (2024 USD):

\begin{itemize}[leftmargin=*]
\item Concrete: $\sim \$0.10$/kg in place.
\item Mild steel structural: $\$0.80$ -- $\$1.50$/kg.
\item Aluminium 6061: $\$3$ -- $\$5$/kg.
\item Stainless 304: $\$3$ -- $\$6$/kg.
\item Stainless 316: $\$5$ -- $\$10$/kg.
\item Brass: $\$8$ -- $\$15$/kg.
\item Titanium Grade 5: $\$30$ -- $\$60$/kg.
\item Carbon-fibre IM tow: $\$30$ -- $\$80$/kg.
\item Aerospace nickel-based superalloy: $\$80$ -- $\$200$/kg.
\item Engineering plastics (PEEK, polyimide): $\$80$ -- $\$200$/kg.
\item Single-crystal silicon (semiconductor grade): $\sim \$500$/kg.
\end{itemize}

The cost is the material price; it does not include the
processing-to-shape cost. A complex titanium part costs ten times more
than its material content because of the machining and the chip waste;
a simple steel part costs little more than its material content.
Process-cost charts (cost per kilogram versus production volume,
parametric by manufacturing route) are the working tool for the
process-cost analysis \cite[ch.~14]{text:ashby2017}.

\engterm{Manufacturability}. The shape and tolerance the design
requires must be achievable by the chosen manufacturing route at the
chosen production volume. Working constraints:

\begin{itemize}[leftmargin=*]
\item Casting requires alloys with good castability (eutectic
  compositions for fluidity, narrow freezing range for shrinkage
  control).
\item Forging requires hot-workable alloys with adequate ductility at
  the forging temperature.
\item Machining requires alloys with reasonable machinability index
  (free-cutting brass, leaded steels; not aerospace titanium).
\item Welding requires weldable alloys (most carbon steels; not
  precipitation-hardened aluminium 7075 without post-weld heat
  treatment).
\item Additive manufacturing requires alloys with adequate fluidity
  for melt-pool stability (Inconel 718, Ti-6Al-4V, 17-4 PH stainless;
  not pure copper without special process control).
\item Composite lay-up requires resin systems and fibre formats
  compatible with the chosen process (autoclave-grade prepreg for
  primary structure; out-of-autoclave grades for cost-bound
  secondary structure).
\end{itemize}

Volume V Chapter 11 develops the manufacturing-process effects on
properties; this chapter notes only that manufacturability is the
discipline-side constraint of the selection step.

\engterm{Availability}. A material chosen in 2024 may not be available
in 2030 in the required quantity, alloy specification, or geographic
delivery time. Working considerations:

\begin{itemize}[leftmargin=*]
\item Strategic materials with restricted geographic supply: rare
  earths (dysprosium, neodymium for permanent magnets), platinum-
  group metals (palladium, rhodium for catalysts), tungsten, cobalt,
  lithium.
\item Materials subject to trade restrictions: titanium grades
  qualified for U.S. aerospace; export-controlled high-performance
  materials.
\item Materials subject to single-source supply: ceramic substrate
  monoliths for automotive catalytic converters; aerospace-grade
  carbon-fibre prepreg.
\end{itemize}

A working design treats availability as a screening constraint: the
material must be available at the required production rate and in the
required specifications, or it cannot be specified. Designs that
depend on a single source are designs at single-source risk.

\engterm{Sustainability}. The material's environmental footprint over
its life cycle includes the embodied energy of production, the
embodied carbon, the toxicity of extraction and processing, the
recyclability of the part at end of life, and the substitution risk
under regulatory tightening. Working figures of embodied carbon (kg
CO$_2$-equivalent per kilogram of material, 2024 values):

\begin{itemize}[leftmargin=*]
\item Concrete (Portland cement basis): $\sim 0.2$ kg CO$_2$/kg
  concrete.
\item Steel (recycled, EAF route): $\sim 0.5$ -- $1.0$.
\item Steel (primary, BOF integrated): $\sim 1.8$ -- $2.5$.
\item Aluminium (recycled): $\sim 0.5$ -- $1.0$.
\item Aluminium (primary): $\sim 8$ -- $20$ (highly grid-dependent;
  hydroelectric primary is at the low end, coal-electricity primary
  at the high end).
\item Carbon-fibre composites: $\sim 20$ -- $40$.
\item Polymers: $\sim 2$ -- $5$ for commodity (PE, PP, PVC), $\sim 10$
  -- $30$ for engineering (PEEK, polyimide).
\end{itemize}

The sustainability constraint is increasingly regulatory: the European
Union's EU Taxonomy and Carbon Border Adjustment Mechanism (CBAM, in
force from 2026 for selected sectors) impose embodied-carbon
disclosures and tariffs on imported goods. The 2030-and-beyond design
cycle treats embodied carbon as a binding constraint rather than an
afterthought.

\engterm{The integrated Ashby chart}. The working selection method
integrates the four non-mechanical constraints by adding them as
screening cuts to the mechanical-property analysis. The shortlist that
survives every screening, including cost, manufacturability,
availability, and sustainability, is the chosen materials set. The
chosen material from the set is the one that maximises the
performance index of the design's objective.

\section{Failure: materials decisions made on price alone}\pathtag{core}

The materials-selection failure pattern that recurs across the
engineering record is price-alone selection: the engineer (or the
purchasing department, or the procurement contract) chose the lowest-
priced material that nominally met the specification, without
weighing the lifecycle cost of the resulting choice. The pattern
appears in several documented forms.

\engterm{Substitution of a nominally-equivalent grade}. A
specification calls for type 316L stainless steel; the supply chain
offers 304L at a lower price; the procurement substitutes 304L on the
argument that ``it is also stainless.'' The substitution is
inadequate for chloride service; the part fails by stress-corrosion
cracking in years rather than decades. The fix is the explicit
specification of the alloy in the contract, with no substitutions
without engineering approval.

\engterm{Substitution of a counterfeit material}. The supply chain
delivers a part stamped with the called-for alloy but, on test, the
chemistry is different. The class of failures includes counterfeit
fasteners in safety-critical structure, counterfeit electronic
components in aerospace and defence, and counterfeit titanium in
aerospace primary structure. The fix is supplier qualification,
material-test reports, and periodic re-test of received lots.

\engterm{The Hubble primary mirror polishing case}. The 1990 Hubble
spherical-aberration failure is sometimes cited as a price-driven
selection failure; the canonical investigation report attributes it
to a measurement error during polishing, not to a materials-choice
error \cite{acc:nasa-hubble-optical-systems-1990}. The relevant
materials-selection lesson is upstream: the ULE titania-silicate
substrate chosen for the mirror was the right material at the cost,
mass, and dimensional-stability constraints; the polishing process
that the substrate's properties enabled was where the failure
occurred. The lesson for the chapter is that a material's nominal
performance can be undone by a process the design failed to
constrain.

\engterm{Price-only selection in cost-bound infrastructure}. Cheap
concrete with high water-to-cement ratio, made and placed without
proper curing, fails by spalling and rebar corrosion within decades
in environments where a slightly more expensive low-permeability mix
would have lasted a century. The aggregate cost of premature
infrastructure replacement (bridge decks, parking garages, marine
piers) across the United States and Europe over $1970$ -- $2020$ is
on the order of trillions of dollars, current as of 2025 estimates;
much of it traces to mid-twentieth-century cost-bound materials
specifications.

\engterm{Price-only selection in safety-critical fasteners}. A high-
strength fastener specification calls for a quenched-and-tempered
grade with a maximum yield strength to control hydrogen-assisted
cracking; the supply chain delivers a grade with higher yield strength
at a lower price; the fastener cracks in service under cathodic-
protection-induced hydrogen embrittlement. The class includes the
hydrogen-cracking incidents in U.K. North Sea oil-platform fasteners
in the 1980s and 1990s that led to revision of the API and DnV
specifications on high-strength-fastener yield-strength limits.

\engterm{The discipline of resisting price-only pressure}. The
engineering response to price-only pressure is a published
specification, an explicit substitution-approval gate, a material-
test-report regime, and a recorded engineering rationale for the
chosen material that survives in the project's documentation. The
discipline is procedural; the failure mode is the failure to enforce
the procedure.

\begin{principle}[Materials selected on price alone fail in service at
predictable mechanisms]
A materials choice made by price without engineering assessment
substitutes the unknown engineering trade-off in the chosen material
for a known cost differential. The unknown trade-off resolves in
service, usually by the failure mode the rejected candidate would
have resisted. The engineering remediation is procedural: a
documented specification, an approved-supplier list, a material-test-
report regime, a substitution-approval gate. Where the procedure is
not enforced, the failure mode is the canonical materials-selection
failure of the engineering record.
\end{principle}

\begin{archetype}[Optimisation and ethics]
A materials-selection decision is an optimisation problem with an
ethics layer. The optimisation chooses among materials by the
explicit performance indices; the ethics layer asks who bears the
cost of a wrong choice. A cost-bound infrastructure choice externalises
the cost of premature failure onto the public; a safety-critical
fastener choice externalises the cost onto the workers and bystanders;
a sustainability-bound choice externalises the cost or the benefit
onto future generations and onto the climate. The chapter's Ashby
method makes the optimisation visible; the ethics is the discipline
of being honest about the externalities that the optimisation does
not see.
\end{archetype}

\begin{estimation}
\textbf{A pedestrian bridge of $20\,\text{m}$ span and $2\,\text{m}$
width is to be designed for minimum embodied carbon, with a $50$-year
service life and a $5\,\text{kN}/\text{m}^2$ live load. Estimate the
relative embodied carbon for the three candidate structural systems:
prestressed concrete, structural steel, and glulam timber.}

\smallskip
\textit{Estimate, before reading on.}

\smallskip
The bending moment at midspan is approximately $M = w L^2 / 8 = 5 \times
2 \times 400 / 8 = 500\,\text{kN}\cdot\text{m}$, with $w = q b = 10\,
\text{kN}/\text{m}$. The required section modulus depends on the
material's design stress: $S = M / \sigma_{\text{allow}}$. Working
allowables: concrete $f_b \approx 10\,\text{MPa}$ (limited by
cracking, not strength); steel $\sigma_{\text{allow}} \approx
160\,\text{MPa}$ (S275 working stress); glulam timber
$\sigma_{\text{allow}} \approx 20\,\text{MPa}$.

The required section moduli are: concrete $S = 0.05\,\text{m}^3$;
steel $S = 3.1 \times 10^{-3}\,\text{m}^3$; glulam $S = 0.025\,
\text{m}^3$.

For a rectangular section $b \times d$ with $b = 2\,\text{m}$, $S =
b d^2 / 6 \Rightarrow d = \sqrt{3 S / b}$. The depths are:
concrete $d \approx 0.27\,\text{m}$; steel (an I-beam, with
flange dominant) $d \approx 0.07\,\text{m}$ effective; glulam $d
\approx 0.19\,\text{m}$.

The volumes (approximately) are: concrete $V \approx 0.27 \times 2
\times 20 = 10.8\,\text{m}^3$ (mass $\sim 26\,\text{t}$); steel
$V \approx 0.5\,\text{m}^3$ (mass $\sim 4\,\text{t}$ for an I-beam
arrangement); glulam $V \approx 7.6\,\text{m}^3$ (mass $\sim
3.4\,\text{t}$).

The embodied carbon is mass times embodied-carbon factor: concrete
$26 \times 0.2 = 5.2\,\text{t CO}_2$; steel $4 \times 1.8 = 7.2\,
\text{t CO}_2$ (primary BOF); glulam $3.4 \times (-0.5) = -1.7\,
\text{t CO}_2$ if sustainably harvested (negative because the wood
sequesters carbon; this depends on the regional forestry practice and
on the end-of-life fate of the timber).

The glulam timber wins on embodied carbon; the steel option is
intermediate; the concrete option is high in mass but the embodied-
carbon factor is low per kilogram of concrete. Recycled-content
steel (EAF route at $\sim 0.5\,\text{t CO}_2$/t) closes the gap with
the concrete option and is competitive with the glulam.

The postmortem: the embodied-carbon calculation is sensitive to the
recycled fraction of the steel, the regional grid mix that powers the
cement kiln, and the assumed end-of-life fate of the timber. A
two-or-threefold uncertainty in the result is honest. The estimate's
value is to show that the three options are within a factor of three
of each other on embodied carbon, not orders of magnitude apart; the
discriminating constraint is likely to be other (durability of the
timber in a humid climate; aesthetics; cost; maintenance burden).
\end{estimation}

\section{Project}\pathtag{core}

\begin{project}[Materials-selection study for a structural component]
The reader specifies a structural component (a bicycle frame, a
backpack frame, a balcony railing, a small footbridge, a pressure
vessel for breathing-gas storage, a kitchen vessel for cooking, a
domestic radiator) and produces an end-to-end materials-selection
report using the Ashby method. The deliverable is a fifteen- to
twenty-five-page report.

\textbf{Track}. Design study with property-chart analysis.
\textbf{Hazard class}: none (analytical project).

\textbf{Procedure}.

\begin{itemize}[leftmargin=*]
\item \engterm{Translation}. Define function, constraints, objective,
  free variables. State each constraint quantitatively (stiffness,
  strength, temperature range, corrosion environment, cost ceiling,
  embodied carbon target).
\item \engterm{Index construction}. Derive the performance index for
  the objective from the loading mode, the geometry, and the binding
  constraint. Show the algebra.
\item \engterm{Chart screening}. Plot the property chart appropriate
  to the index (modulus-density, strength-density, or another),
  apply the constraint screening cuts, and identify the candidate
  shortlist by sweeping the constant-index line.
\item \engterm{Documentation}. For each shortlist candidate, look up
  the working numbers from Callister, Ashby's catalogue, or
  manufacturer data sheets. Tabulate.
\item \engterm{Multi-objective}. If the design has two or more
  objectives, plot the Pareto front and identify the chosen material
  with an explicit utility statement.
\item \engterm{Non-mechanical constraints}. Apply cost, manufacturability,
  availability, and sustainability as screening or as secondary
  objectives. Document the rejections.
\item \engterm{Final recommendation}. State the chosen material with
  the reasoning, the principal alternative, and the conditions
  under which the alternative would be preferred.
\end{itemize}

\textbf{Reflection ($2{,}000$ to $3{,}000$ words)}. The reader
addresses:

\begin{itemize}[leftmargin=*]
\item What was the binding constraint of the design (the constraint
  that, if relaxed, would change the chosen material)?
\item How sensitive is the choice to assumed values of cost, embodied
  carbon, or the utility weights?
\item Which engineering failure modes would the chosen material be
  most likely to exhibit, and what inspection or design provisions
  are required?
\item One feature of the comparison that surprised the reader.
\end{itemize}

\textbf{Deliverable}. The translation document, the index derivation,
the property charts (annotated by hand or with software), the
candidate-shortlist table with sources, the multi-objective analysis
(if applicable), the final recommendation memo, and a one-paragraph
summary suitable for a non-engineering client.

\textbf{Time}. Six to twelve hours of analysis, four to eight hours of
chart construction and annotation, four to six hours of write-up.
\end{project}

\section{Exercises}

\subsection*{Calculation}\pathtag{core}

\begin{exercise}
A tie rod of length $L$ must carry a tensile force $F$. Derive the
mass of the rod in terms of $L$, $F$, the design stress $\sigma$, and
the density $\rho$. Identify the performance index for minimum mass.
\end{exercise}

\begin{exercise}
A simply-supported beam of length $L$ and rectangular cross-section
$b \times d$ (both free) must carry a midspan force $F$ at a midspan
deflection $\delta$. Derive the performance index for minimum mass.
\end{exercise}

\begin{exercise}
Compute the performance index $M = E^{1/2} / \rho$ for the
following materials and rank them: aluminium 6061 ($E = 70\,\text{GPa}$,
$\rho = 2700\,\text{kg}/\text{m}^3$); titanium Grade 5 ($E = 114\,
\text{GPa}$, $\rho = 4430$); carbon-fibre composite ($E = 140\,\text{GPa}$,
$\rho = 1600$); structural steel ($E = 200\,\text{GPa}$, $\rho = 7850$);
balsa wood ($E = 4\,\text{GPa}$, $\rho = 200$).
\end{exercise}

\begin{exercise}
Compute the embodied-carbon-weighted cost per unit performance
$\rho C_{\text{CO}_2} / M$ for the same five materials, using the
embodied-carbon figures in the text.
\end{exercise}

\begin{exercise}
A pressure vessel of radius $R$ and wall thickness $t$ is to be
designed for a working pressure $P$ at a hoop stress $\sigma_y$. Derive
the mass of the vessel in terms of $R$, $P$, $\sigma_y$, and $\rho$.
Identify the performance index for minimum mass.
\end{exercise}

\begin{exercise}
A thermal-storage element must hold a thermal energy $Q$ at a
temperature $T_{\max}$. Derive the volume required as a function of
$\rho$, $c_p$, $T_{\max} - T_{\text{ambient}}$, and the latent-heat
contribution if a phase-change material is used. Identify the performance
index for minimum volume.
\end{exercise}

\subsection*{Derivation}\pathtag{standard}

\begin{exercise}
Derive the performance index for a light, stiff plate (bending,
stiffness objective, free thickness, two in-plane dimensions fixed),
$M = E^{1/3} / \rho$.
\end{exercise}

\begin{exercise}
Derive the performance index for a damage-tolerant light tie,
$M = K_{Ic} / \rho$, for a part with a tolerable through-crack of fixed
length.
\end{exercise}

\begin{exercise}
Derive the performance index for a heat-exchanger tube wall of fixed
diameter and free thickness, with the objective of minimum mass per
unit heat-transfer rate.
\end{exercise}

\begin{exercise}
Show that the performance index for a torsion spring (storing maximum
elastic energy per unit mass) is $\sigma_y^2 / (E \rho)$.
\end{exercise}

\subsection*{Estimation}\pathtag{core}

\begin{exercise}
Estimate the mass of a light, stiff bicycle frame using the
$E^{1/2}/\rho$ index and a target tube stiffness. Compare aluminium
6061 and carbon-fibre composite tubes.
\end{exercise}

\begin{exercise}
Estimate the embodied carbon of a typical mid-rise apartment building
($10$ floors, $40 \times 20\,\text{m}$ footprint) constructed in
reinforced concrete versus structural steel versus mass timber.
\end{exercise}

\begin{exercise}
Estimate the price differential between mild steel and stainless 316
for a $5{,}000\,\text{kg}$ industrial vessel, and the corresponding
break-even service life assuming the stainless option needs no
maintenance and the mild-steel option requires annual repainting at
$\$1{,}000$.
\end{exercise}

\subsection*{Simulation}\pathtag{mastery}

\begin{exercise}
Implement a script that reads a CSV of material properties ($E$,
$\sigma_y$, $\rho$, cost per kg, embodied carbon per kg) and computes
the rank of materials for a chosen performance index. Plot the chosen
index against price for the top 10 candidates.
\end{exercise}

\begin{exercise}
Implement a multi-objective Pareto-front identifier: given a set of
materials and two objective functions, return the Pareto-optimal
subset. Verify on a published Ashby chart problem.
\end{exercise}

\subsection*{Design}\pathtag{standard}

\begin{exercise}
Specify a material for a low-cost domestic radiator ($60\,\degC$
operating temperature, $10$-year service life, embodied carbon as
secondary objective). State the chosen material, the alternative, and
the rationale.
\end{exercise}

\begin{exercise}
Specify a material for an outdoor-mounted electrical-enclosure box
($-30\,\degC$ to $60\,\degC$ service, UV exposure, low electrical
conductivity, IP65 rating). State the chosen material, the principal
alternative, and the rejected candidates.
\end{exercise}

\begin{exercise}
Specify a material for a single-piece kitchen-knife blade
(hardness, edge retention, corrosion resistance in food-contact
environment, ease of sharpening). State the chosen material, the
alternative, and the rationale.
\end{exercise}

\begin{exercise}
Specify a material for the load-bearing structure of a small electric
delivery vehicle (mass, cost, manufacturability, recyclability at
end of life). State the chosen material and the trade-off analysis.
\end{exercise}

\subsection*{Diagnosis and reverse engineering}\pathtag{core}

\begin{exercise}
A pressure vessel chosen for a chemical-process service is found to
be made of mild steel rather than the specified stainless-316. The
service environment is hot ($90\,\degC$) chloride-containing water.
Diagnose the substitution's likely consequence and recommend the
inspection programme until the vessel can be replaced.
\end{exercise}

\begin{exercise}
A consumer product (a cooking pan) has been delivered with the
nominally same alloy as the prototype but a different surface
treatment. The pan develops surface staining and food-contact
contamination within months. Diagnose the difference and the
materials-selection-process step that failed.
\end{exercise}

\begin{exercise}
A wind-turbine blade designed for $20$-year service was retired at
$8$ years because of glass-fibre fatigue cracking. Diagnose the
materials-selection step where the design margin against fatigue was
inadequate and the data the engineer would now require.
\end{exercise}

\subsection*{Failure analysis}\pathtag{core}

\begin{exercise}
A municipal water-main rupture is traced to galvanic-coupling
corrosion between an aluminium fitting and a brass valve. Identify
the materials-selection error, the standard remediation, and the
mechanism by which the rupture occurred.
\end{exercise}

\begin{exercise}
A bridge-deck reinforcement-corrosion problem is traced to a
late-design substitution of epoxy-coated rebar for galvanised rebar
on cost grounds. The epoxy coating's installation tolerances were
not enforced. Diagnose the materials-selection-process step that
failed.
\end{exercise}

\subsection*{Judgment}\pathtag{standard}

\begin{exercise}
A start-up's pitch claims a new alloy that outperforms aerospace-grade
titanium at one-fifth the cost. Write a one-paragraph reply identifying
the questions about strength, fatigue, processing route, and
qualification regime that the next investor meeting must answer.
\end{exercise}

\begin{exercise}
A facility manager proposes to standardise on a single ``low-cost,
high-strength'' carbon steel for all structural service across a
chemical plant, on procurement-simplification grounds. Write a one-
paragraph evaluation, identifying the corrosion-service categories
the proposal would harm.
\end{exercise}

\begin{exercise}
A regulator proposes to require an embodied-carbon disclosure for
every primary structural material specified on a public-works
project, with no other constraints changed. Identify the predictable
unintended consequences and the procedural safeguards required to
prevent them.
\end{exercise}

\begin{exercise}
A colleague proposes to choose a material by computing a
multi-objective utility function with weights derived from a market
survey of customer preferences. Write a one-paragraph reply on the
limits of customer preferences as a guide to safety-critical
materials selection.
\end{exercise}
